%% Complete proofs for the subdirect representation toolkit.

\section{Subdirect products and irreducible quotients}

\begin{annexitem}{annex:p10.subdirect.1}{thm:p10-subdirect-kernels}{The congruence-intersection criterion}
\begin{proof}
First, an intersection of equivalence relations is an equivalence
relation, and compatibility with a basic operation holds in an
intersection because it holds in each constituent congruence.
For the empty intersection the universal relation is a congruence.
These facts justify the meets, joins, and generated congruences used
in the body: common upper bounds always include \(\nabla_A\).

For a homomorphism \(h:A\to B\), its kernel is an equivalence relation.  If
\(h(a_j)=h(b_j)\) for \(1\leq j\leq n\), then
\[
 h(f^A(a_1,\ldots,a_n))
  =f^B(h(a_1),\ldots,h(a_n))
  =f^B(h(b_1),\ldots,h(b_n))
  =h(f^A(b_1,\ldots,b_n)),
\]
so it is a congruence.  The nullary case is just equality of the
image of the same constant with itself.
The induced quotient operations make \(q_\theta\) a surjective
homomorphism with kernel \(\theta\).

For the product map \(e\), the homomorphism identity holds at every
coordinate because it holds for each \(h_i\); coordinate equality
then proves it in the product.  This also covers constant symbols.
Moreover,
\[
 e(a)=e(b)
 \quad\Longleftrightarrow\quad
 \forall i\in I\ h_i(a)=h_i(b)
 \quad\Longleftrightarrow\quad
 (a,b)\in\bigcap_{i\in I}\ker h_i.
\]
Thus \(e\) is injective exactly when that intersection is
\(\Delta_A\).  Its image is closed under every basic operation,
and the inverse of an injective homomorphism onto its image
preserves operations: apply the homomorphism identity to the
preimages of an arbitrary tuple in the image.
So injectivity gives an isomorphism onto that subalgebra.
If each \(h_i\) is surjective, the coordinate maps on the image
are surjective as well, proving the subdirect assertion.

For \(\theta=\bigcap_i\theta_i\), define
\[
 \bar e([a]_\theta)=([a]_{\theta_i})_{i\in I}.
\]
This is well-defined since \(\theta\subseteq\theta_i\) for every
\(i\).  It is a homomorphism by the quotient and product operation
clauses.  Equality of two image tuples says
\((a,b)\in\theta_i\) for every \(i\), which is exactly
\([a]_\theta=[b]_\theta\), proving injectivity.
Every coordinate is onto: the class \([a]_{\theta_i}\) is the
\(i\)-coordinate of \(\bar e([a]_\theta)\).

If \(I=\varnothing\), the product contains the unique empty tuple,
the intersection is \(\nabla_A\), and \(A/\nabla_A\) has one
element.  The resulting map is the unique isomorphism between
these one-element algebras; coordinate surjectivity is vacuous.
The map with domain \(A\) itself is injective in this case exactly
when \(\Delta_A=\nabla_A\), namely when \(A\) is trivial.
\end{proof}
\end{annexitem}

\begin{annexitem}{annex:p10.subdirect.2}{thm:p10-congruence-correspondence}{Congruences above a quotient kernel}
\begin{proof}
Suppose \(\theta\subseteq\phi\in\operatorname{Con}A\).
If \([a]_\theta=[a']_\theta\) and
\([b]_\theta=[b']_\theta\), then
\((a,a'),(b,b')\in\theta\subseteq\phi\).
Symmetry and transitivity of \(\phi\) therefore give
\[
 (a,b)\in\phi\quad\Longleftrightarrow\quad(a',b')\in\phi.
\]
Hence the relation \(\phi/\theta\) is independent of representatives.
Its equivalence-relation laws follow from those of \(\phi\).
For any basic \(n\)-ary operation, choose representatives of the
finitely many argument classes.  Compatibility of \(\phi\) and
the definition of quotient operations prove compatibility of
\(\phi/\theta\).  Constants cause no additional condition.
Thus \(\phi/\theta\in\operatorname{Con}(A/\theta)\).

Conversely, the inverse image of a congruence \(\psi\) on \(A/\theta\)
under the homomorphism \(q_\theta\) is an equivalence relation
compatible with every operation.  It contains \(\theta\), since
equal quotient classes are \(\psi\)-related.
These two assignments are inverse:
\[
 ([a]_\theta,[b]_\theta)\in\phi/\theta
       \quad\Longleftrightarrow\quad(a,b)\in\phi,
\]
and every quotient class has a representative.
They preserve and reflect inclusion by the same equivalence.
An order isomorphism between complete lattices carries a least
upper bound to a least upper bound, and a greatest lower bound
to a greatest lower bound, including empty bounds.
Finally \(\theta/\theta=\Delta_{A/\theta}\).
\end{proof}
\end{annexitem}

\begin{annexitem}{annex:p10.subdirect.3}{thm:p10-subdirect-monolith}{Characterizing subdirect irreducibility}
\begin{proof}
Suppose first that \(A\) is subdirectly irreducible and
\((\theta_i)_{i\in I}\) is a nonempty family with intersection
\(\Delta_A\).
Theorem~\ref{thm:p10-subdirect-kernels} gives a subdirect embedding
\(A\to\prod_i A/\theta_i\).
Some quotient coordinate must be an isomorphism, and its kernel
is then \(\Delta_A\).  Since that kernel is \(\theta_i\), the
family contains equality.

Conversely, suppose the stated intersection property holds, and
let \(e:A\to\prod_i A_i\) be a subdirect embedding with nonempty
index set.  The kernels of \(\pi_i\circ e\) intersect in
\(\Delta_A\).
One kernel is equality, so the corresponding coordinate map is
injective as well as surjective.  A bijective homomorphism of
algebras is an isomorphism by the inverse-operation argument in
Annex~\ref{annex:p10.subdirect.1}.
This proves subdirect irreducibility.

If \(A\) is trivial, all its congruences equal \(\Delta_A\);
every nonempty congruence family consequently has an equality
member.  Equivalently, a surjection from a one-element algebra
has a one-element codomain, and every coordinate of a subdirect
embedding into a nonempty family is an isomorphism.
The nonempty-index hypothesis supplies a coordinate to choose.

Now suppose \(A\) is nontrivial.  The set
\[
 J=\operatorname{Con}A\setminus\{\Delta_A\}
\]
is nonempty because it contains \(\nabla_A\).
If \(A\) is subdirectly irreducible, then
\(\bigcap J\ne\Delta_A\): otherwise the family of all members
of \(J\) would violate the intersection property just proved.
As \(\bigcap J\) is itself a congruence, it is then a member of
\(J\) below every other member, so it is a monolith.
Conversely, if \(J\) has a least member \(\mu\), its intersection
is \(\mu\ne\Delta_A\).
For any nonempty family of congruences all strictly above equality,
the intersection contains \(\mu\); it therefore cannot be equality.
This establishes the intersection property and all equivalences.

If \(a\ne b\) and \((a,b)\in\mu\), the principal congruence
\(\operatorname{Cg}_A(a,b)\) is strictly above equality and is
contained in \(\mu\).  Leastness of \(\mu\) gives the reverse
inclusion, proving \(\operatorname{Cg}_A(a,b)=\mu\).
\end{proof}
\end{annexitem}

\begin{annexitem}{annex:p10.subdirect.4}{thm:p10-subdirect-representation}{Birkhoff's subdirect representation theorem}
\begin{proof}
Suppose first that \(A\) has at least two elements.
Fix distinct \(a,b\in A\), and consider the partially ordered set
\[
 \mathcal P_{a,b}
   =\{\theta\in\operatorname{Con}A:(a,b)\notin\theta\},
 \qquad\text{ordered by inclusion}.
\]
It is a set and is nonempty, since \(\Delta_A\in\mathcal P_{a,b}\).

\paragraph{Upper bounds for chains.}
For a nonempty chain \(\mathcal D\subseteq\mathcal P_{a,b}\),
put \(\lambda=\bigcup_{\theta\in\mathcal D}\theta\).
It is reflexive because each member of the chain contains
\(\Delta_A\), and symmetric because each member is symmetric.
If \((x,y)\in\theta_1\) and \((y,z)\in\theta_2\), one of
\(\theta_1,\theta_2\) contains the other; its transitivity gives
\((x,z)\in\lambda\).

For an \(n\)-ary basic operation \(f\) with \(n>0\), suppose
\((x_j,y_j)\in\lambda\) for \(1\leq j\leq n\).
There are only finitely many pairs.  Select a chain member
containing each pair, and take the largest of those finitely
many chain members.  It contains all the pairs, and its
compatibility with \(f\) puts
\[
 (f^A(x_1,\ldots,x_n),f^A(y_1,\ldots,y_n))\in\lambda.
\]
For a nullary operation compatibility follows from reflexivity.
Thus \(\lambda\) is a congruence.  This step uses finite arity,
not finiteness of the signature or the carrier.
It still omits \((a,b)\), since membership in the union would
put that pair in a member of \(\mathcal D\).
Hence \(\lambda\in\mathcal P_{a,b}\) is a chain upper bound.
The empty chain has the upper bound \(\Delta_A\); its set-theoretic
union need not be used.

Zorn's lemma now gives a maximal member \(\theta_{a,b}\) of
\(\mathcal P_{a,b}\).
This is the choice principle used in the argument.

\paragraph{The separating quotient is irreducible.}
Put
\[
 \eta_{a,b}
       =\theta_{a,b}\vee\operatorname{Cg}_A(a,b).
\]
It strictly contains \(\theta_{a,b}\), since it contains
\((a,b)\).
If \(\phi\) is any congruence strictly containing
\(\theta_{a,b}\), its omission of \((a,b)\) would contradict
maximality in \(\mathcal P_{a,b}\).
It therefore contains \((a,b)\), hence contains
\(\operatorname{Cg}_A(a,b)\) as well as \(\theta_{a,b}\).
By the definition of join,
\(\eta_{a,b}\subseteq\phi\).
Thus \(\eta_{a,b}\) is the least congruence strictly above
\(\theta_{a,b}\).

The correspondence of
Theorem~\ref{thm:p10-congruence-correspondence} sends this least
member to
\(\eta_{a,b}/\theta_{a,b}\), the least congruence strictly above
equality on \(A/\theta_{a,b}\).
The quotient is nontrivial, since \([a]_{\theta_{a,b}}\ne
[b]_{\theta_{a,b}}\).
Theorem~\ref{thm:p10-subdirect-monolith} therefore proves that
\(A/\theta_{a,b}\) is subdirectly irreducible.

\paragraph{Combining the quotients.}
Use the axiom of choice to select such a maximal congruence for
each member of the set
\[
 P=\{(a,b)\in A^2:a\ne b\}.
\]
Every pair in \(P\) is omitted by its selected congruence, while
every diagonal pair belongs to every congruence.  Consequently
\[
 \bigcap_{(a,b)\in P}\theta_{a,b}=\Delta_A.
\]
The congruence-intersection criterion gives a subdirect embedding
\[
 A\longrightarrow\prod_{(a,b)\in P}A/\theta_{a,b}.
\]
Each coordinate is the surjective quotient homomorphism, so all
factors are homomorphic images of \(A\), as required.
The embedding identifies \(A\) with its image subalgebra, not
necessarily with the entire product.

If \(A\) is trivial, it is subdirectly irreducible by the preceding
theorem, and the identity gives the one-factor representation.
Alternatively \(P\) is empty and the unique map from \(A\) to
the empty product is an isomorphism of one-element algebras.

If \(A\) is finite and nontrivial, \(P\) has
\(|A|(|A|-1)\) elements and every factor is finite.
Each \(\mathcal P_{a,b}\) is a nonempty finite partially ordered
set, so it has a maximal member without Zorn's lemma.
Only finitely many selections are needed; finite choice is
provable without the axiom of choice.
The same intersection and quotient arguments give the stated
finite representation and bound.
\end{proof}
\end{annexitem}
